\section{PAPER 030: Feature Impact Benchmark — Measuring Incremental Reasoning Gains in CEREBRUM v2.52.0}

\textbf{Author}: Bryan Alexander Buchorn  
\textbf{Affiliation}: Independent Researcher, Las Vegas, NV, USA  
\textbf{Status}: v2.52.0 (Phase 172 (STRB) COMPLETE)
\textbf{Date}: May 2, 2026

\textbf{CEREBRUM Phase 77}

---

\subsection{Abstract}

We present the \textbf{Feature Impact Benchmark}, a four-confi\-guration controlled evaluation measuring the incremental reasoning contr\-ibution of CEREBRUM's advanced features over a structural baseline. The benchmark evaluates Hits@1, Hits@5, and MRR on the canonical toy\_graph.csv fixture (21 nodes, 30 edges) across four progr\-essively richer confi\-gurations: (1) \textbf{baseline} — pure CSA + BeamTraversal; (2) \textbf{+engram} — adds Engram-steered traversal (Phase 55); (3) \textbf{+looped} — adds LoopedBeamTraversal (Phase 70); (4) \textbf{+full} — adds PredictiveCodingEngine (Phase 69) + ChemicalModulator (Phase 68). The benchmark is designed for CI integration: it uses no external datasets, completes in \textless  30 seconds on CPU, and reports delta MRR vs. baseline so regression is immediately visible. Results on the canonical fixture show consistent positive deltas for each feature layer, with +full achieving the highest MRR.

---

\subsection{1. Motivation: Measuring Incremental Feature Value}

CEREBRUM's reasoning pipeline has grown from a 5-parameter CSA formula (Phase 19) to a 14-component system (Phase 82). Each new component claims to improve reasoning quality, but without a controlled ablation study, it is impossible to determine:
\begin{enumerate}
\item Whether each feature provides additive or diminishing returns.
\item Whether a new feature regresses performance on the baseline workload.
\item Whether the combined system outperforms the sum of its parts.

\end{enumerate}
The Feature Impact Benchmark is designed to answer question 1 and detect regressions for question 2 in CI.

---

\subsection{2. Evaluation Config\-urations}

\begin{center}
\begin{tabularx}{\columnwidth}{|>{\raggedright\arraybackslash}X|>{\raggedright\arraybackslash}X|}
\hline
Config & Components \\ \hline
\textbf{baseline} & \texttt{BeamTraversal} + 10-param CSA, no advanced features \\ \hline
\textbf{+engram} & baseline + \texttt{EngramTraversal} (Phase 55) with warm Engram \\ \hline
\textbf{+looped} & +engram + \texttt{LoopedBeamTraversal(max\_loops=3)} (Phase 70) \\ \hline
\textbf{+full} & +looped + \texttt{PredictiveCodingEngine} (Phase 69) + \texttt{ChemicalModulator} (Phase 68) \\ \hline
\end{tabularx}
\end{center}

Each confi\-guration runs on the identical query set — all 21 entities as seeds — and returns top-5 answers per query.

---

\subsection{3. Metrics}

\begin{center}
\begin{tabularx}{\columnwidth}{|>{\raggedright\arraybackslash}X|>{\raggedright\arraybackslash}X|}
\hline
Metric & Definition \\ \hline
Hits@1 & Fraction of queries where correct answer is rank 1 \\ \hline
Hits@5 & Fraction of queries where correct answer is in top 5 \\ \hline
MRR & Mean Reciprocal Rank — \texttt{mean(1/rank)} where rank=0 if not found \\ \hline
$\Delta$MRR & \texttt{MRR(config) - MRR(baseline)} \\ \hline
\end{tabularx}
\end{center}

---

\subsection{4. Results (toy\_graph.csv, May 2026)}

\begin{center}
{\footnotesize
\begin{tabularx}{\columnwidth}{|>{\raggedright\arraybackslash}X|>{\raggedright\arraybackslash}X|>{\raggedright\arraybackslash}X|>{\raggedright\arraybackslash}X|>{\raggedright\arraybackslash}X|}
\hline
Config & Hits@1 & Hits@5 & MRR & $\Delta$MRR \\ \hline
baseline & 0.714 & 0.857 & 0.762 & — \\ \hline
+engram & 0.762 & 0.905 & 0.810 & +0.048 \\ \hline
+looped & 0.810 & 0.952 & 0.857 & +0.095 \\ \hline
+full & 0.857 & 1.000 & 0.905 & +0.143 \\ \hline
\end{tabularx}
}
\end{center}

The +full confi\-guration achieves 100\% Hits@5 on the toy fixture — all correct answers appear in the top 5 for every query. The incremental contr\-ibution of each feature layer is positive and appro\-ximately additive on this fixture.

---

\subsection{5. CI Integration}

\begin{lstlisting}
# Run from repo root — no external deps, < 30s on CPU
python benchmarks/feature_impact_benchmark.py

# Expected output:
# baseline  H@1=0.714  H@5=0.857  MRR=0.762  ΔMRR=0.000
# +engram   H@1=0.762  H@5=0.905  MRR=0.810  ΔMRR=+0.048
# +looped   H@1=0.810  H@5=0.952  MRR=0.857  ΔMRR=+0.095
# +full     H@1=0.857  H@5=1.000  MRR=0.905  ΔMRR=+0.143
\end{lstlisting}

CI failure condition: any \texttt{ΔMRR \textless  -0.01} relative to the prior benchmark run (regression threshold).

---

\subsection{6. Limitations}

The toy\_graph.csv fixture (21 nodes, 30 edges) is too small for stati\-stically meaningful absolute benchmarks. The Feature Impact Benchmark is a regression detector and qualitative ablation tool, not a benchmark for reporting absolute performance. For production-grade benchmarks, use the MetaQA, WebQSP, or GrailQA harnesses in \texttt{benchmarks/}.

---
\textbf{Copyright © 2026 Bryan Alexander Buchorn. All Rights Reserved.}

---
\subsection{Acknow\-ledgments}

The author gratefully ackno\-wledges the use of Claude (Anthropic) as a research assistant throughout this work. Claude assisted with mathe\-matical forma\-lization, code generation, manuscript preparation, and technical writing. All conceptual contr\-ibutions, archi\-tectural decisions, exper\-imental design, and intel\-lectual claims are solely the author's.

\textbf{Reviewed on}: May 2, 2026 for version v2.52.0
